\subsection{一个集合到另一个集合的映射集}

设 E, F 为集合。从 E 到 F 的映射的图是 \(E \times F\) 的子集。因此， \(\mathfrak{P}(E \times F)\) 中由从 E 到 F 的映射的图所构成的元素集合，是 \(\mathfrak{P}(E \times F)\) 的一个子集，记为 \(F^{E}\) 。所有三元组 \(f = (G, E, F)\) （其中 \(G \in F^{E}\) ）所成的集合，就是从 E 到 F 的映射全体；它记为 \(\mathcal{F}(E, F)\) 。显然， \(G \to (G, E, F)\) 是从 \(F^{E}\) 到 \(\mathcal{F}(E, F)\) 的一个双射（称为典范双射）。这个双射使我们能够立即把关于集合 \(F^{E}\) 的每个命题转化为关于 \(\mathcal{F}(E, F)\) 的命题，反之亦然。

设 E、E'、F、F' 为集合。设 u 为从 E' 到 E 的映射，v 为从 F 到 F' 的映射。则函数 \(f \rightarrow v \circ f \circ u\) 是从 \(\mathcal{F}(\mathrm{E}, \mathrm{F})\) 到 \(\mathcal{F}(\mathrm{E}', \mathrm{F}')\) 的映射。

命题 2. (1) 若 \(u\) 是从 \(\mathbf{E}'\) 到 \(\mathbf{E}\) 的满射，而 \(v\) 是从 \(\mathbf{F}\) 到 \(\mathbf{F}'\) 的单射，则映射 \(f \to v \circ f \circ u\) 是单射。

\begin{enumerate}
\def\labelenumi{(\arabic{enumi})}
\setcounter{enumi}{1}
\tightlist
\item
  如果 \(u\) 是从 \(\mathbf{E}'\) 到 \(\mathbf{E}\) 的单射，而 \(v\) 是从 \(\mathbf{F}\) 到 \(\mathbf{F}'\) 的满射，则映射 \(f \to v \circ f \circ u\) 是满射。
\end{enumerate}

我们将假设集合 E、E'、F、F' 均为非空，否则命题显然成立。

\begin{enumerate}
\def\labelenumi{(\arabic{enumi})}
\tightlist
\item
  让 \(s\) 是……的一个截面 \(u\) 并令 \(r\) 为 \(v\) （§3，定义 11）。则 \(r \circ (v \circ f \circ u) \circ s = I_F \circ f \circ I_E = f\) ，因此
\end{enumerate}

\[
f \rightarrow v \circ f \circ u
\]

是单射的。

\begin{enumerate}
\def\labelenumi{(\arabic{enumi})}
\setcounter{enumi}{1}
\tightlist
\item
  让 \(r'\) 为 \(u\) 并令 \(s'\) 是……的一个截面 \(v\) 。对于每个映射 \(f': E' \to F'\) 我们拥有 \(v \circ (s' \circ f' \circ r') \circ u = f'\) ，这表明映射 \(f \to v \circ f \circ u\) 的满射。
\end{enumerate}

推论. 如果 \(u\) 是从 \(\mathbf{E}'\) 到 \(\mathbf{E}\) 的双射，且 \(v\) 是从 \(\mathbf{F}\) 到 \(\mathbf{F}'\) 的双射，那么 \(f \to v \circ f \circ u\) 是双射。

设 A、B、C 为三个集合，并设 f 为从 \(B \times C\) 到 A 的一个映射。对每一个 \(y \in C\) ，令 \(f(\bullet, y)\) 为从 B 到 A 的偏映射 \(x \to f(x, y)\) （§ 3，第 9 段）；函数 \(y \to f(\bullet, y)\) 是从 C 到 \(\mathcal{F}(B, A)\) 的映射。反之，对每个从 C 到 \(\mathcal{F}(B, A)\) 的映射 g，都存在唯一一个映射 \(f\) （从 \(\mathbf{B} \times \mathbf{C}\) 到 \(\mathbf{A}\) ），使得 \(g(y) = f(\bullet, y)\) （对每个 \(y \in \mathbf{C}\) ），它就是映射 \((x, y) \to (g(y))(x)\) 。因此：

命题3. 对每个映射 \(f\) （从 \(\mathbf{B} \times \mathbf{C}\) 到 \(\mathbf{A}\) ），以 \(\tilde{f}\) 表示映射 \(y \to f(\bullet, y)\) （从 \(\mathbf{C}\) 到 \(\mathcal{F}(\mathbf{B}, \mathbf{A})\) ），则函数 \(f \to \tilde{f}\) 是从 \(\mathcal{F}(\mathbf{B} \times \mathbf{C}, \mathbf{A})\) 到 \(\mathcal{F}(\mathbf{C}, \mathcal{F}(\mathbf{B}, \mathbf{A}))\) 的双射（称为典范双射）。

同样地，可定义从 \(\mathcal{F}(\mathbf{B}\times\mathbf{C},\mathbf{A})\) 到 \(\mathcal{F}(\mathbf{B},\mathcal{F}(\mathbf{C},\mathbf{A}))\) 的一个典范双射。由映射与函数图之间的一一对应，这些双射又给出从 \(\mathbf{A}^{\mathbf{B}\times\mathbf{C}}\) 到 \((\mathbf{A}^{\mathbf{B}})^{\mathbf{C}}\) （相应地， \((\mathbf{A}^{\mathbf{C}})^{\mathbf{B}})\) ）的典范双射。

